\begingroup
\renewcommand{\arraystretch}{1.45}
\fancyhead[L]{\small ARITHMETISCH-GEOMETRISCHES MITTEL}
\Needspace{36\baselineskip}
\AGMArticle{5}

\textsc{Aufgabe.}\quad \textit{Das arithmetisch-geometrische Mittel zwischen der Zahl $1+x$, die größer als eins ist, und eins durch eine nach Potenzen von $x$ fortschreitende Reihe auszudrücken.}

\textit{Lösung.}\quad Da $\mathrm M(1,1)=1$ ist, nehmen wir an
\[
\mathrm M(1+x,1)=1+h'x+h''x^2+h'''x^3+h''''x^4+\text{usw.}
\]
wobei $h'$, $h''$, $h'''$, $h''''$ konstante, von $x$ unabhängige Koeffizienten sind.
Es sei $x=2t+tt$; dann ist
\[
\begin{aligned}&\mathrm M(1+x,1)=\mathrm M(1+t+\tfrac12tt,\ 1+t)\\&=(1+t)\mathrm M(1+\frac{\tfrac12tt}{1+t},1).\end{aligned}
\]
Daher ergibt sich
\[
\begin{aligned}
&1+h'(2t+tt)+h''(2t+tt)^2+h'''(2t+tt)^3+\text{usw.}\\[6pt]
={}&1+t+h'(\tfrac12tt)+h''\frac{\tfrac14t^4}{1+t}+h'''\frac{\tfrac18t^6}{(1+t)^2}+\text{usw.}
\end{aligned}
\]
\Needspace{29\baselineskip}
Hieraus ergeben sich die Gleichungen
\[
\begin{aligned}
2h'&=1\\[6pt]
4h''+h'&=\tfrac12h'\\[6pt]
8h'''+4h''&=0\\[6pt]
16h''''+12h'''+h''&=\tfrac14h''\\[6pt]
32h^{\mathrm V}+32h''''+6h'''&=-\tfrac14h''\\[6pt]
64h^{\mathrm{VI}}+80h^{\mathrm V}+24h''''+h'''&=\tfrac14h''+\tfrac18h'''\\[6pt]
128h^{\mathrm{VII}}+192h^{\mathrm{VI}}+80h^{\mathrm V}+8h''''&=-\tfrac14h''-\tfrac28h'''\\[6pt]
256h^{\mathrm{VIII}}+448h^{\mathrm{VII}}+240h^{\mathrm{VI}}+40h^{\mathrm V}+h''''&=\tfrac14h''+\tfrac38h'''+\tfrac1{16}h''''
\end{aligned}
\]
usw.\quad woraus folgt
\[
\begin{gathered}h'=\tfrac12,\quad h''=-\tfrac1{16},\quad h'''=\tfrac1{32},\\ h''''=-\tfrac{21}{1024},\quad h^{\mathrm V}=\tfrac{31}{2048},\quad h^{\mathrm{VI}}=-\tfrac{195}{16384}\end{gathered}
\]
Daher
\[
\begin{aligned}&\mathrm M(1+x,1)=1+\tfrac12x-\tfrac1{16}xx+\tfrac1{32}x^3\\&\quad-\tfrac{21}{1024}x^4+\tfrac{31}{2048}x^5-\tfrac{195}{16384}x^6\ \text{usw.}\end{aligned}
\]
Übrigens ist leicht zu erkennen, dass das Mittel zwischen $1$ und der Zahl
\SourceLeaf{375}
\(1-x\), die kleiner als eins ist, \(1-h'x+h''xx-h'''x^3+\text{usw.}=1-\frac12x-\frac1{16}xx-\frac1{32}x^3-\frac{21}{1024}x^4-\text{usw.}\) sein wird.
Da diese Koeffizienten kein offenkundiges Gesetz erkennen lassen, gehen wir über diese Reihen hinweg
und versuchen einen anderen Weg, der zu größerem Erfolg führen wird.

\Needspace{29\baselineskip}
\AGMArticle{6}

\textsc{Aufgabe.} \textit{Das arithm.-geom. Mittel zwischen \(1+x\) und \(1-x\) durch eine nach Potenzen von \(x\) fortschreitende Reihe auszudrücken.}

\textit{Lösung.} Da
\[
\mathrm{M}(1+x,1-x)=(1-x)\mathrm{M}(1+\frac{2x}{1-x},1)
\]
ergibt sich aus der Reihe des vorigen Artikels sofort, wenn dort \(\frac{2x}{1-x}\) für \(x\) eingesetzt wird,
\[
\begin{aligned}&\mathrm{M}(1+\frac{2x}{1-x},1)=1+x+\frac34xx+\frac34x^3\\&\quad+\frac{43}{64}x^4+\frac{43}{64}x^5+\frac{161}{256}x^6+\text{usw.}\end{aligned}
\]
und hieraus
\[
\mathrm{M}(1+x,1-x)=1-\frac14xx-\frac5{64}x^4-\frac{11}{256}x^6+\text{usw.}
\]

\Needspace{31\baselineskip}
Die Koeffizienten dieser Reihe lassen sich auch unabhängig von der Reihe des vorigen Artikels durch das folgende Verfahren ermitteln. Es werde \(x=\nobreak{}\frac{2t}{1+tt}\) gesetzt; dann ist
\[
\begin{aligned}&\mathrm{M}(1+x,1-x)=\mathrm{M}(\frac{1-tt}{1+tt},1)\\&=\frac1{1+tt}\mathrm{M}(1+tt,1-tt)\end{aligned}
\]
Daher ergibt sich, wenn
\[
\mathrm{M}(1+x,1-x)=1+\alpha xx+\beta x^4+\gamma x^6+\delta x^8+\text{usw.}
\]
gesetzt wird (denn dass Potenzen von \(x\) mit ungeradem Exponenten nicht vorkommen, ist von selbst klar)
\[
\begin{aligned}
&(1+tt)\{1+\alpha(\frac{2t}{1+tt})^2+\beta(\frac{2t}{1+tt})^4\\&\quad+\gamma(\frac{2t}{1+tt})^6+\delta(\frac{2t}{1+tt})^8+\text{usw.}\}\\[6pt]
&\qquad=1+\alpha t^4+\beta t^8+\gamma t^{12}+\delta t^{16}+\text{usw.}
\end{aligned}
\]
\Needspace{29\baselineskip}
Hieraus ergeben sich die Gleichungen
\[
\begin{aligned}
1+4\alpha&=0\\&\qquad\text{daraus }\alpha=-\frac14\\[6pt]
-4\alpha+16\beta&=\alpha\\&\qquad\beta=-\frac5{64}\\[6pt]
4\alpha-48\beta+64\gamma&=0\\&\qquad\gamma=-\frac{11}{256}\\[6pt]
-4\alpha+96\beta-320\gamma+256\delta&=\beta\\&\qquad\delta=-\frac{29}{1024}-\frac5{16384}=-\frac{469}{16384}\\[6pt]
&\text{usw.}\qquad\text{usw.}
\end{aligned}
\]
\Needspace{23\baselineskip}
Auch in dieser Reihe folgen die Koeffizienten keinem einfachen Gesetz; wird aber eins durch
diese Reihe dividiert, so ergibt sich
\SourceLeaf{376}
\[
\begin{aligned}&\frac1{\mathrm{M}(1+x,1-x)}=1+\frac14xx+\frac9{64}x^4+\frac{25}{256}x^6\\&\quad+\frac{1225}{16384}x^8+\text{usw.}\end{aligned}
\]
wobei wir auf den ersten Blick sehen, dass die Koeffizienten die Quadrate der Wurzeln
\(\frac12,\allowbreak \ \frac12.\frac34,\allowbreak \ \frac12.\frac34.\frac56,\allowbreak \ \frac12.\frac34.\frac56.\frac78\)
sind und somit nach einem sehr einfachen Gesetz fortschreiten. Doch das Verfahren, durch
das wir zu diesem sehr schönen Schluss gelangt sind, hat lediglich die Kraft
eines induktiven Schlusses; wir nehmen uns vor, ihn in den folgenden Untersuchungen
zur Gewissheit zu erheben.

\clearpage
\endgroup
